\begin{abstract}
Receiver-dependent signal transformations complicate closed-set radio frequency fingerprinting when a classifier is transferred to unseen physical receivers. We consolidate a leakage-audited WiSig evaluation that distinguishes source-only generalization, bounded unlabeled receiver adaptation, and learned receiver-context conditioning. A six-transmitter subset is evaluated by leave-one-receiver-out protocols over 32 physical receivers, using five fixed seeds and disjoint support and query pools. Target-adaptive methods receive the same 128 unlabeled receiver packets; attentive context conditioning uses 32 peers per query. T3A prototype adjustment achieves receiver-equal macro-F1 of \NTThreeAMacroFOne{}, compared with \NPZeroMacroFOne{} for source-only empirical risk minimization and \NPTwoMacroFOne{} for attentive conditioning. The T3A improvement is \TGain{}, positive on 31 receivers, with a receiver-bootstrap interval of [\TLower{}, \TUpper{}]. The later benchmark's receiver sign-flip analysis also supports a positive difference, but was specified after earlier WiSig findings and is not an independent confirmation. Attention-based conditioning is sensitive to receiver-matched support yet provides no meaningful average advantage over the independent baseline. Benefits depend strongly on available support, and lower classification error does not guarantee lower expected calibration error. These results motivate information-matched receiver-adaptation evaluation, not a new superior context architecture. Evidence is limited to one dataset and deterministically partitioned support, not physically acquired calibration episodes.
\end{abstract}
